\subsection{Diagnostic: learning rate $10^{-4}$}
To test whether optimisation explains the gap, we repeated strategy~2 with a single change, lr
$10^{-3}\to10^{-4}$ (same split, seed, resolution, epochs and snapshots). Training became stable (no collapses;
training loss 0.27--0.32 after 20 epochs versus 0.61 at the end of the paper-setting run) and every
metric improved (Table~\ref{tab:main}): QWK \SIIlowBestQWK{} for the best-validation model (epoch 11) and
\SIIlowVoteQWK{} for the 10-snapshot vote, accuracy \SIIlowVoteAcc{}, macro-F1 \SIIlowBestMF{}.
Minority recall (mean of Severe and PDR) rose from \SIIBestMinRec{} to \SIIlowBestMinRec{}: Severe 0.00$\to$0.31
and PDR 0.30$\to$0.55 for the single model, and predictions off by two or more grades halved (48$\to$24).
With stable training the snapshots are uniformly good (test QWK 0.834--0.875), and the vote now matches the
best single model, consistent with the snapshot-ensemble assumption that members are of comparable quality.

This closes about half of the gap to 0.954 (remaining $\approx$0.07 QWK) and indicates that the paper's stated
learning rate, at least at our batch size, is the main cause of the faithful run's poor result. Severe NPDR
remains the weakest grade (single-model recall 0.31; 0.10 for the vote), and Mild is still mostly confused with
Moderate (24/56). Leakage again has no measurable effect (QWK 0.887 with and without the 29 duplicated test
images). This run is the baseline we will carry into Stage~4.
